\documentclass[../../main.tex]{subfiles}

\begin{document}

\section{完备度量空间}\label{section:ch07-complete-metric}

在这一节中, 我们定义完备性的概念, 并且证明若 $Y$ 是完备度量空间, 则函数空间 $\mathcal{C}(X,Y)$ 在一致度量下也是完备的. 此外还要证明每一个度量空间能被等距地嵌入到一个完备度量空间中.

\begin{definition}\label{definition:ch07-cauchy-complete}
设 $(X,d)$ 为度量空间. $X$ 的点的序列 $(x_n)$ 称为 $(X,d)$ 中的一个 Cauchy 序列(Cauchy sequence), 如果它具有以下性质: 任意给定 $\varepsilon>0$, 存在一个整数 $N$, 使得
\begin{align*}
d(x_n,x_m)<\varepsilon
\end{align*}
对任意 $n,m\geqslant N$ 成立. 度量空间 $(X,d)$ 称为完备的(complete), 若 $X$ 中的每一个 Cauchy 序列收敛.
\end{definition}

当然, $X$ 中任何一个收敛序列必定是一个 Cauchy 序列, 完备性是要求其逆命题也成立.

注意, 完备度量空间 $(X,d)$ 的闭子集 $A$, 在限制度量下必定是完备的. 因为 $A$ 中 Cauchy 序列也是 $X$ 中的 Cauchy 序列, 所以在 $X$ 中收敛; 又因为 $A$ 是 $X$ 的闭子集, 所以其极限必定在 $A$ 中.

另外, 如果 $X$ 关于度量 $d$ 是完备的, 那么 $X$ 在相应于 $d$ 的标准有界度量
\begin{align*}
\bar{d}(x,y)=\min\{d(x,y),1\}
\end{align*}
下也是完备的, 反之也对. 因为一个序列 $(x_n)$ 对于 $\bar{d}$ 是 Cauchy 序列当且仅当它对于 $d$ 是 Cauchy 序列, 并且一个序列对于 $\bar{d}$ 收敛当且仅当它对于 $d$ 收敛.

以下引理是可以用来判别度量空间是否完备的一个有用的准则.

\begin{lemma}\label{lemma:ch07-cauchy-subsequence}
如果度量空间 $X$ 中每一个 Cauchy 序列有一个收敛的子序列, 那么 $X$ 是完备的.
\end{lemma}

\begin{proof}
设 $(x_n)$ 是 $(X,d)$ 中的一个 Cauchy 序列. 我们证明, 如果 $(x_n)$ 有一个子序列 $(x_{n_i})$ 收敛于一点 $x$, 则序列 $(x_n)$ 本身也收敛于 $x$.

给定 $\varepsilon>0$, 首先选取足够大的 $N$, 使得对于所有的 $n,m\geqslant N$, 有
\begin{align*}
d(x_n,x_m)<\varepsilon/2.
\end{align*}
(因为 $(x_n)$ 是一个 Cauchy 序列.) 然后选取足够大的整数 $i$, 使 $n_i\geqslant N$ 并且
\begin{align*}
d(x_{n_i},x)<\varepsilon/2.
\end{align*}
(因为 $n_1<n_2<\cdots$ 是一个递增的整数序列, 并且 $(x_{n_i})$ 收敛于 $x$.) 综合上述事实, 就可以得到所要证的结果: 对于 $n\geqslant N$,
\begin{align*}
d(x_n,x)\leqslant d(x_n,x_{n_i})+d(x_{n_i},x)<\varepsilon.
\end{align*}
\end{proof}

\begin{theorem}\label{theorem:ch07-euclidean-complete}
对于通常度量, 即欧氏度量 $d$ 和平方度量 $\rho$, 欧氏空间 $\mathbb{R}^k$ 都是完备的.
\end{theorem}

\begin{proof}
为了证明度量空间 $(\mathbb{R}^k,\rho)$ 是完备的, 令 $(x_n)$ 为 $(\mathbb{R}^k,\rho)$ 中的一个 Cauchy 序列. 集合 $\{x_n\}$ 是 $(\mathbb{R}^k,\rho)$ 的有界子集. 因为如果选择 $N$, 使得对于所有的 $n,m\geqslant N$, 有
\begin{align*}
\rho(x_n,x_m)\leqslant 1,
\end{align*}
那么数
\begin{align*}
M=\max\{\rho(x_1,\mathbf{0}),\cdots,\rho(x_{N-1},\mathbf{0}),\rho(x_N,\mathbf{0})+1\}
\end{align*}
是 $\rho(x_n,\mathbf{0})$ 的一个上界. 于是, 序列 $(x_n)$ 中的点都在立方体 $[-M,M]^k$ 中. 因为这个立方体是紧致的, 根据第 28 节中的定理, 序列 $(x_n)$ 有一个收敛子序列, 于是 $(\mathbb{R}^k,\rho)$ 是完备的.

为了证明 $(\mathbb{R}^k,d)$ 是完备的, 注意一个序列关于 $d$ 是 Cauchy 序列当且仅当它关于 $\rho$ 是 Cauchy 序列, 并且一个序列相对于 $d$ 收敛当且仅当它相对于 $\rho$ 收敛.
\end{proof}

现在我们来处理积空间 $\mathbb{R}^\omega$ 的问题. 为此我们需要下面这个关于积空间中的序列的引理.

\begin{lemma}\label{lemma:ch07-product-sequence}
设 $X=\prod X_\alpha$ 是积空间, $x_n$ 是 $X$ 的一个点的序列. 则 $x_n\to x$ 当且仅当对于每一个 $\alpha$, $\pi_\alpha(x_n)\to\pi_\alpha(x)$.
\end{lemma}

\begin{proof}
这一结论曾在第 19 节中作为习题给出, 在此我们给出它的证明. 因为投射 $\pi_\alpha\colon X\to X_\alpha$ 是连续的, 所以收敛序列像还是收敛序列, 所以命题的"必要性"成立. 下面来证明充分性. 设对于每一个 $\alpha\in J$, 有 $\pi_\alpha(x_n)\to\pi_\alpha(x)$. 设 $U=\prod U_\alpha$ 是 $X$ 的包含 $x$ 的一个基元素. 对于 $\alpha$, 如果 $U_\alpha$ 不等于整个空间 $X_\alpha$, 选取 $N_\alpha$ 使得当 $n\geqslant N_\alpha$ 时 $\pi_\alpha(x_n)\in U_\alpha$. 设 $N$ 是 $N_\alpha$ 中最大的, 那么当 $n\geqslant N$ 时 $x_n\in U$.
\end{proof}

\begin{theorem}\label{theorem:ch07-Romega-complete}
对于积空间 $\mathbb{R}^\omega$, 总存在一个度量, 使得相对这个度量而言 $\mathbb{R}^\omega$ 是完备的.
\end{theorem}

\begin{proof}
设 $\bar{d}(a,b)=\min\{|a-b|,1\}$ 是 $\mathbb{R}$ 的标准有界度量. 设 $D$ 是 $\mathbb{R}^\omega$ 的一个度量, 定义为
\begin{align*}
D(\mathbf{x},\mathbf{y})=\sup\{\bar{d}(x_i,y_i)/i\}.
\end{align*}
那么 $D$ 诱导了 $\mathbb{R}^\omega$ 的积拓扑. 下面将证明 $\mathbb{R}^\omega$ 关于 $D$ 是完备的. 设 $\mathbf{x}_n$ 是 $(\mathbb{R}^\omega,D)$ 的一个 Cauchy 序列. 由于
\begin{align*}
\bar{d}(\pi_i(\mathbf{x}),\pi_i(\mathbf{y}))\leqslant iD(\mathbf{x},\mathbf{y}),
\end{align*}
对于固定的 $i$, 序列 $\pi_i(\mathbf{x}_n)$ 是 $\mathbb{R}$ 中的一个 Cauchy 序列, 因此它收敛, 不妨设其收敛于 $a_i$. 于是序列 $\mathbf{x}_n$ 收敛到 $\mathbb{R}^\omega$ 中的点 $\mathbf{a}=\{a_1,a_2,\cdots\}$.
\end{proof}

\begin{example}\label{example:ch07-Q-not-complete}
一个非完备的度量空间的例子就是关于通常度量 $d(x,y)=|x-y|$ 的有理数空间 $\mathbb{Q}$. 例如, (在 $\mathbb{R}$ 中)收敛于 $\sqrt{2}$ 的有限小数序列
\begin{align*}
1.4,1.41,1.414,1.4142,1.41421,\cdots
\end{align*}
是 $\mathbb{Q}$ 中的一个 Cauchy 序列, 但 (在 $\mathbb{Q}$ 中)不收敛.
\end{example}

\begin{example}\label{example:ch07-open-interval-not-complete}
另一个非完备的空间是关于度量 $d(x,y)=|x-y|$ 的 $\mathbb{R}$ 中的开区间 $(-1,1)$. 在这个空间中, 定义为
\begin{align*}
x_n=1-1/n
\end{align*}
的序列 $(x_n)$ 是一个 Cauchy 序列, 但不收敛. 这个例子表明, 完备性不是拓扑性质. 也就是说, 它不是在同胚下保持不变的性质. 因为 $(-1,1)$ 同胚于实直线 $\mathbb{R}$, 而 $\mathbb{R}$ 在通常度量下是完备的.
\end{example}

虽然积空间 $\mathbb{R}^n$ 和 $\mathbb{R}^\omega$ 都有度量使之成为完备空间, 但对于积空间 $\mathbb{R}^J$ 来说, 未必能有这样的结果. 因为如果 $J$ 不可数, 则 $\mathbb{R}^J$ 甚至不能度量化(见第 21 节). 但是, 在集合 $\mathbb{R}^J$ 上存在不同的拓扑, 其一便是由一致度量诱导的. 我们将看到, 在这一度量下, $\mathbb{R}^J$ 是完备的.

我们定义一致度量如下:

\begin{definition}\label{definition:ch07-uniform-metric}
设 $(Y,d)$ 是一个度量空间, $\bar{d}(a,b)=\min\{d(a,b),1\}$ 是 $Y$ 上相应于 $d$ 的标准有界度量. 若 $x=(x_\alpha)_{\alpha\in J}$ 与 $y=(y_\alpha)_{\alpha\in J}$ 是笛卡儿积 $Y^J$ 的两个点, 令
\begin{align*}
\bar{\rho}(\mathbf{x},\mathbf{y})=\sup\{\bar{d}(x_\alpha,y_\alpha)\mid\alpha\in J\}.
\end{align*}
容易验证, $\bar{\rho}$ 是一个度量. 称其为 $Y^J$ 的相应于 $Y$ 的度量 $d$ 的一致度量(uniform metric).
\end{definition}

我们在这里对 $Y^J$ 中的元素使用了标准"串"记号. 因为 $Y^J$ 中的元素是从 $J$ 到 $Y$ 的函数, 我们也可以使用函数记号表示它们. 在这一章中, 函数记号会比串记号更方便些, 因此我们全部使用函数记号. 在此定义中一致度量的定义具有以下形式: 若 $f,g\colon J\to Y$, 则
\begin{align*}
\bar{\rho}(f,g)=\sup\{\bar{d}(f(\alpha),g(\alpha))\mid\alpha\in J\}.
\end{align*}

\begin{theorem}\label{theorem:ch07-uniform-power-complete}
如果空间 $Y$ 关于度量 $d$ 是完备的, 则空间 $Y^J$ 在相应于 $d$ 的一致度量 $\bar{\rho}$ 下也是完备的.
\end{theorem}

\begin{proof}
我们曾经说过, 如果 $(Y,d)$ 是完备的, 那么 $(Y,\bar{d})$ 也是完备的, 其中 $\bar{d}$ 是相应于 $d$ 的有界度量. 现在假定 $f_1,f_2,\cdots$ 是 $Y^J$ 中点的一个序列, 它相对于 $\bar{\rho}$ 是一个 Cauchy 序列. 给定 $\alpha\in J$, 对于所有 $m$ 和 $n$ 有
\begin{align*}
\bar{d}(f_n(\alpha),f_m(\alpha))\leqslant\bar{\rho}(f_n,f_m),
\end{align*}
这说明 $f_1(\alpha),f_2(\alpha),\cdots$ 是 $(Y,\bar{d})$ 中的一个 Cauchy 序列. 因此, 这个序列收敛, 设它收敛于点 $y_\alpha$. 设函数 $f\colon J\to Y$ 定义为 $f(\alpha)=y_\alpha$. 我们断言: 序列 $(f_n)$ 关于度量 $\bar{\rho}$ 收敛于 $f$.

给定 $\varepsilon>0$, 首先选取充分大的 $N$, 使得 $\bar{\rho}(f_n,f_m)<\varepsilon/2$, 其中 $n,m\geqslant N$. 于是, 特别是对于 $n,m\geqslant N$ 及 $\alpha\in J$,
\begin{align*}
\bar{d}(f_n(\alpha),f_m(\alpha))<\varepsilon/2.
\end{align*}
固定 $n$ 和 $\alpha$, 让 $m$ 变得任意大, 则只要 $n\geqslant N$, 对任意 $\alpha\in J$ 有
\begin{align*}
\bar{d}(f_n(\alpha),f(\alpha))\leqslant\varepsilon/2.
\end{align*}
因此, 当 $n\geqslant N$ 时, 有
\begin{align*}
\bar{\rho}(f_n,f)\leqslant\varepsilon/2<\varepsilon.
\end{align*}
定理证毕.
\end{proof}

现在, 我们可以考虑集合 $Y^X$, 其中 $X$ 是一个拓扑空间, 而不单单是一个集合了. 当然, 这并不会对前面所作的讨论有所影响. 当我们考虑所有函数 $f\colon X\to Y$ 构成的集合时, $X$ 的拓扑是无所谓的. 但是, 假定我们考虑的是由所有连续函数 $f\colon X\to Y$ 构成的 $Y^X$ 的子集 $\mathcal{C}(X,Y)$ 时, 就与 $X$ 的拓扑有关了. 这时, 我们得出, 若 $Y$ 是完备的, 则在一致度量下 $\mathcal{C}(X,Y)$ 也是完备的. 这一结论对于所有有界函数 $f\colon X\to Y$ 构成的集合 $\mathcal{B}(X,Y)$ 也成立. (一个函数 $f$ 称为有界的(bounded), 若它的像集 $f(X)$ 是度量空间 $(Y,d)$ 的一个有界子集.)

\begin{theorem}\label{theorem:ch07-C-B-closed}
设 $X$ 是一个拓扑空间, $(Y,d)$ 是一个度量空间. 则连续函数集 $\mathcal{C}(X,Y)$ 和有界函数集 $\mathcal{B}(X,Y)$ 在一致度量下都是 $Y^X$ 中的闭集. 因此, 如果 $Y$ 是完备的, 那么这两个空间在一致度量下都是完备的.
\end{theorem}

\begin{proof}
这个定理的前半部分就是一致极限定理的一个翻版. 首先, 我们证明, 如果 $Y^X$ 中的序列 $(f_n)$ 相对于 $Y^X$ 中的度量 $\bar{\rho}$ 收敛于 $Y^X$ 的元素 $f$, 那么在第 21 节中所给的定义下相对于 $Y$ 上的度量 $\bar{d}$, 它也一致收敛于 $f$. 因为, 给定 $\varepsilon>0$, 选取整数 $N$, 使得对于所有 $n>N$,
\begin{align*}
\bar{\rho}(f,f_n)<\varepsilon.
\end{align*}
于是, 对于所有 $x\in X$ 及所有的 $n\geqslant N$, 有
\begin{align*}
\bar{d}(f_n(x),f(x))\leqslant\bar{\rho}(f_n,f)<\varepsilon.
\end{align*}
所以 $(f_n)$ 一致收敛于 $f$.

现在, 我们证明 $\mathcal{C}(X,Y)$ 相对于度量 $\bar{\rho}$ 是 $Y^X$ 的闭集. 设 $f$ 为 $Y^X$ 的一个元素, 同时也是 $\mathcal{C}(X,Y)$ 的一个极限点. 于是存在 $\mathcal{C}(X,Y)$ 中一个元素序列 $(f_n)$ 关于度量 $\bar{\rho}$ 收敛于 $f$. 根据一致极限定理, $f$ 连续, 所以 $f$ 属于 $\mathcal{C}(X,Y)$.

最后, 我们要证明 $\mathcal{B}(X,Y)$ 是 $Y^X$ 的一个闭集. $f$ 是 $\mathcal{B}(X,Y)$ 的一个极限点, 存在 $\mathcal{B}(X,Y)$ 中一个元素序列 $(f_n)$ 收敛于 $f$. 选取 $N$ 足够大, 使得 $\bar{\rho}(f_N,f)<1/2$. 从而对于 $x\in X$, 我们有 $\bar{d}(f_N(x),f(x))<1/2$. 由此可见 $d(f_N(x),f(x))<1/2$. 所以, 如果 $f_N(X)$ 的直径为 $M$, 那么 $f(X)$ 的直径最多为 $M+1$. 因此, $f\in\mathcal{B}(X,Y)$.

综上所述, 我们有: 若 $Y$ 关于度量 $d$ 是完备的, 则 $\mathcal{C}(X,Y)$ 和 $\mathcal{B}(X,Y)$ 在度量 $\bar{\rho}$ 下都是完备的.
\end{proof}

\begin{definition}\label{definition:ch07-sup-metric}
设 $(Y,d)$ 是一个度量空间. 可以在由 $X$ 到 $Y$ 的有界函数构成的集合 $\mathcal{B}(X,Y)$ 上定义另一度量如下:
\begin{align*}
\rho(f,g)=\sup\{d(f(x),g(x))\mid x\in X\}.
\end{align*}
易见 $\rho$ 的定义是确切的, 因为当 $f(X)$ 和 $g(X)$ 都有界时, 集合 $f(X)\cup g(X)$ 也是有界的. 这个度量就被称为上确界度量(sup metric).
\end{definition}

在上确界度量和一致度量之间有一些简单的联系. 事实上, 若 $f,g\in\mathcal{B}(X,Y)$, 则有
\begin{align*}
\bar{\rho}(f,g)=\min\{\rho(f,g),1\}.
\end{align*}
因为如果 $\rho(f,g)>1$, 那么至少存在一点 $x_0\in X$, 有 $d(f(x_0),g(x_0))>1$, 因此 $\bar{d}(f(x_0),g(x_0))=1$, 并且 $\bar{\rho}(f,g)=1$. 反之, 如果 $\rho(f,g)\leqslant 1$, 那么对任意 $x\in X$, 有 $\bar{d}(f(x),g(x))=d(f(x),g(x))\leqslant 1$, 从而 $\bar{\rho}(f,g)=\rho(f,g)$. 因此在 $\mathcal{B}(X,Y)$ 上, 度量 $\bar{\rho}$ 正是相对于度量 $\rho$ 的标准有界度量. 这就是我们引用记号 $\bar{\rho}$ 表示一致度量的原因, 参见第 20 节.

如果 $X$ 是紧致的, 那么每一个连续函数 $f\colon X\to Y$ 都是有界的. 因此上确界度量就定义在 $\mathcal{C}(X,Y)$ 上. 如果 $Y$ 在度量 $d$ 下是完备的, 那么 $\mathcal{C}(X,Y)$ 在相应的一致度量 $\bar{\rho}$ 下也是完备的, 因此在上确界度量 $\rho$ 下也是完备的. 在这种情况下, 我们常使用上确界度量而不是一致度量.

现在我们要证明一个经典定理, 得到的结论是: 每一个度量空间都可以等距地嵌入到一个完备度量空间中. (还有一种证明方式更为直接些, 参见习题 9.) 尽管我们并不会用到这个定理, 但它在数学的其他方面很有用.

\begin{theorem}\label{theorem:ch07-isometric-embedding}
设 $(X,d)$ 是一个度量空间. 则存在一个从 $X$ 到某个完备度量空间中的等距嵌入.
\end{theorem}

\begin{proof}
设 $\mathcal{B}(X,\mathbb{R})$ 是将 $X$ 映入 $\mathbb{R}$ 中的所有有界函数构成的集合. 令 $x_0$ 为 $X$ 中取定的一个点. 给定 $a\in X$, 定义 $\phi_a\colon X\to\mathbb{R}$ 为
\begin{align*}
\phi_a(x)=d(x,a)-d(x,x_0).
\end{align*}
我们断言 $\phi_a$ 是有界的. 因为根据不等式
\begin{align*}
d(x,a)\leqslant d(x,b)+d(a,b),\qquad d(x,b)\leqslant d(x,a)+d(a,b),
\end{align*}
可见
\begin{align*}
|d(x,a)-d(x,b)|\leqslant d(a,b).
\end{align*}
令 $b=x_0$, 我们得出结论, 对于所有 $x\in X$ 有 $|\phi_a(x)|\leqslant d(a,x_0)$.

定义 $\Phi\colon X\to\mathcal{B}(X,\mathbb{R})$ 为
\begin{align*}
\Phi(a)=\phi_a.
\end{align*}
我们证明 $\Phi$ 是从 $(X,d)$ 到完备度量空间 $(\mathcal{B}(X,\mathbb{R}),\rho)$ 中的一个等距嵌入, 也就是说, 证明对于点的任意偶对 $a,b\in X$, 有
\begin{align*}
\rho(\phi_a,\phi_b)=d(a,b).
\end{align*}

根据定义
\begin{align*}
\rho(\phi_a,\phi_b)&=\sup\{|\phi_a(x)-\phi_b(x)|;x\in X\}\\
&=\sup\{|d(x,a)-d(x,b)|;x\in X\},
\end{align*}
可见
\begin{align*}
\rho(\phi_a,\phi_b)\leqslant d(a,b).
\end{align*}
另一方面, 这个不等式可能是不严格的. 因为当 $x=a$ 时, 有
\begin{align*}
|d(x,a)-d(x,b)|=d(a,b).
\end{align*}
\end{proof}

\begin{definition}\label{definition:ch07-completion}
设 $X$ 是一个度量空间. 如果 $h\colon X\to Y$ 是从 $X$ 到完备度量空间 $Y$ 中的一个等距嵌入, 则 $Y$ 的子空间 $\overline{h(X)}$ 是一个完备度量空间. 它称为 $X$ 的一个完备化(completion).
\end{definition}

$X$ 的完备化在等距的意义上说是唯一确定的. 见习题 10.

\subsection*{习题}

\begin{enumerate}
\item 设 $X$ 是一个度量空间.
\begin{enumerate}[(a)]
\item 假定对某一个 $\varepsilon>0$, $X$ 中的每一个 $\varepsilon$-球都有紧致闭包. 证明 $X$ 是完备的.
\item 假定对于每一个 $x\in X$, 存在 $\varepsilon>0$, 使得球 $B(x,\varepsilon)$ 有紧致闭包. 举例说明, $X$ 未必是完备的.
\end{enumerate}
\item 设 $(X,d_X)$ 和 $(Y,d_Y)$ 都是度量空间, $Y$ 是完备的. 设 $A\subset X$. 证明: 如果 $f\colon A\to Y$ 一致连续, 则 $f$ 可以唯一地扩张成一个连续函数 $g\colon\overline{A}\to Y$, 并且 $g$ 是一致连续的.
\item $X$ 的两个度量 $d$ 和 $d'$ 称为度量等价的(metrically equivalent), 如果恒等映射 $i\colon (X,d)\to(X,d')$ 及其逆映射都是一致连续的.
\begin{enumerate}[(a)]
\item 证明: $d$ 度量等价于相对于 $d$ 的标准有界度量 $\bar{d}$.
\item 证明: 若 $d$ 和 $d'$ 是度量等价的, 则 $X$ 在度量 $d$ 下是完备的, 当且仅当它关于度量 $d'$ 是完备的.
\end{enumerate}
\item 证明: 度量空间 $(X,d)$ 是完备的, 当且仅当对 $X$ 中满足 $|A_n|\to 0$ 的任意一个非空闭集套序列 $A_1\supset A_2\supset\cdots$, 有 $\bigcap A_n\neq\varnothing$.
\item 设 $(X,d)$ 是一个度量空间. 以前说过, 一个连续映射 $f\colon X\to X$ 称为一个压缩映射, 指的是如果存在一个数 $\alpha<1$, 使得对于所有的 $x,y\in X$, 有
\begin{align*}
d(f(x),f(y))\leqslant\alpha d(x,y).
\end{align*}
证明: 如果 $f$ 是一个完备度量空间上的压缩映射, 则存在唯一的一个点 $x\in X$, 使得 $f(x)=x$. 与第 28 节中的习题 7 比较.
\item 空间 $X$ 称为拓扑完备的(topologically complete), 如果存在着诱导出 $X$ 的拓扑的一个度量, 使得在这一度量下 $X$ 是完备的.
\begin{enumerate}[(a)]
\item 证明: 一个拓扑完备空间的闭子空间也是拓扑完备的.
\item 证明: 拓扑完备空间的可数积空间 (关于积拓扑) 仍然是拓扑完备的.
\item 证明: 拓扑完备空间的开子空间是拓扑完备的. [提示: 如果 $U\subset X$, 并且 $X$ 关于度量 $d$ 是完备的, 定义 $\phi\colon U\to\mathbb{R}$ 为
\begin{align*}
\phi(x)=1/d(x,X-U).
\end{align*}
通过 $f(x)=x\times\phi(x)$ 将 $U$ 嵌入到 $X\times\mathbb{R}$ 中.]
\item 证明: 如果 $A$ 是拓扑完备空间中的一个 $G_\delta$ 集, 则 $A$ 是拓扑完备的. [提示: 设 $A$ 是开集 $U_n$ 的交, $n\in\mathbb{Z}_+$, 考虑 $A$ 到 $\prod U_n$ 的对角线嵌入 $f(a)=(a,a,\cdots)$.] 证明: 无理数集是拓扑完备的.
\end{enumerate}
\item 证明所有使得 $\sum x_i^2$ 收敛的序列 $(x_1,x_2,\cdots)$ 的集合在 $\ell^2$-度量下是完备的 (参见第 20 节习题 8).
\item 设 $X$ 和 $Y$ 是两个空间, 定义
\begin{align*}
e\colon X\times\mathcal{C}(X,Y)\longrightarrow Y
\end{align*}
为 $e(x,f)=f(x)$, 映射 $e$ 称为一个赋值映射(evaluation map). 证明: 如果 $d$ 是 $Y$ 上的度量, $\mathcal{C}(X,Y)$ 有相应的一致拓扑, 那么 $e$ 是连续的. 在第 46 节中我们将推广这一结论.
\item 设 $(X,d)$ 是一个度量空间. 证明: 存在从 $X$ 到完备度量空间 $(Y,D)$ 的一个等距嵌入 $h$ 如下: 令 $\widetilde{X}$ 表示 $X$ 中点的所有 Cauchy 序列
\begin{align*}
\mathbf{x}=(x_1,x_2,\cdots)
\end{align*}
的集合. 若
\begin{align*}
d(x_n,y_n)\longrightarrow 0.
\end{align*}
则定义 $\mathbf{x}\sim\mathbf{y}$. 用 $[\mathbf{x}]$ 表示 $\mathbf{x}$ 的等价类, $Y$ 表示这些等价类构成的集合. 定义 $Y$ 上的度量 $D$ 为
\begin{align*}
D([\mathbf{x}],[\mathbf{y}])=\lim_{n\to\infty}d(x_n,y_n).
\end{align*}
\begin{enumerate}[(a)]
\item 证明: $\sim$ 是一个等价关系, $D$ 是一个定义确切的度量.
\item 定义 $h\colon X\to Y$ 为 $h(x)$ 是常值序列 $(x,x,\cdots)$ 的等价类:
\begin{align*}
h(x)=[(x,x,\cdots)].
\end{align*}
证明 $h$ 是一个等距嵌入.
\item 证明 $h(X)$ 在 $Y$ 中稠密. 事实上, 给定 $\mathbf{x}=(x_1,x_2,\cdots)\in\widetilde{X}$, 证明 $Y$ 中点的序列 $h(x_n)$ 收敛于 $Y$ 中的点 $[\mathbf{x}]$.
\item 证明: 如果 $A$ 是度量空间 $(Z,\rho)$ 的一个稠密子集, 并且 $A$ 中任意一个 Cauchy 序列在 $Z$ 中收敛, 则 $Z$ 是完备的.
\item 证明 $(Y,D)$ 是完备的.
\end{enumerate}
\item 定理(完备化的唯一性) 设 $h\colon X\to Y$ 和 $h'\colon X\to Y'$ 分别是从度量空间 $(X,d)$ 到完备度量空间 $(Y,D)$ 和 $(Y',D')$ 的两个等距嵌入. 则有一个 $(\overline{h(X)},D)$ 与 $(\overline{h'(X)},D')$ 之间的等距同构在子空间 $h(X)$ 上等于 $h'h^{-1}$.
\end{enumerate}

\end{document}